\chapter{応用：潜在空間の最適輸送幾何}
\label{ch:latent}

本章では，深層生成モデルの潜在空間に Wasserstein 距離で幾何を入れる
Roy--Hauberg
（\emph{Optimal Latent Transport}, NeurIPS 2022 Workshop on
Symmetry and Geometry in Neural Representations）の枠組みを，
本資料の道具で定式化する．
第\ref{ch:wasserstein-metrics}章--第\ref{ch:concentration}章の結果が
どこで使われるかを明示し，同論文が証明なしに用いる主張には
証明または出典を与える．

\section{生成モデルの測度論的定式化}
\label{sec:lt-setup}

潜在空間を $Z=\R^d$（Euclid 距離．
注意~\ref{rem:tc-rn-polish}より Polish 空間），
データ空間を可分距離空間 $(\mathcal{X},d_{\mathcal{X}})$ とする．

\begin{definition}{デコーダ・エンコーダ・集約事後分布}{lt-vae}
  \textbf{デコーダ}とは，$Z$ から $\mathcal{X}$ への Markov 核
  $K$（定理~\ref{thm:meas-disintegration}の (i)(ii)）のことをいい，
  $K(z,\cdot)$ を潜在点 $z$ の\textbf{デコーダ分布}とよぶ．
  prior $\gamma^{(d)}$（標準 Gaussian 測度：
  定義~\ref{def:tc-gaussian}）に対し，混合測度
  （主張~\ref{clm:meas-kernel-mixture}）
  \[
    p_G:=K_*\gamma^{(d)}\in\Prob(\mathcal{X})
  \]
  を\textbf{生成分布}という．
  同様に，\textbf{エンコーダ}とはデータ空間から潜在空間への
  Markov 核 $Q$ のことをいい，データ分布
  $P\in\Prob(\mathcal{X})$ に対する混合測度
  \[
    q_Z:=Q_*P\in\Prob(Z)
  \]
  を\textbf{集約事後分布}という．
\end{definition}

\section{ELBO の正則化項と潜在空間の輸送距離}
\label{sec:lt-elbo}

変分オートエンコーダ（VAE）は，再構成項と正則化項
$\int_{\mathcal{X}}H_{\gamma^{(d)}}\bigl(Q(x,\cdot)\bigr)\,\d P(x)$
（各データ点の事後分布と prior の相対エントロピーの平均）からなる
目的関数（ELBO）を最適化する．
この正則化項が集約事後分布と prior の $W_2$ 距離を制御することを示す．

\begin{lemma}{混合はエントロピーを増やさない}{lt-entropy-convex}
  $\Theta,\mathcal{X}$ を可分距離空間，
  $P\in\Prob(\Theta)$，$\nu\in\Prob(\mathcal{X})$ とし，
  可測関数 $\rho\colon\Theta\times\mathcal{X}\to[0,\infty)$ は
  各 $\theta\in\Theta$ で
  $\int_{\mathcal{X}}\rho(\theta,x)\,\d\nu(x)=1$ をみたすとする．
  \[
    \bar\rho(x):=\int_\Theta\rho(\theta,x)\,\d P(\theta)
  \]
  とおくと，$\bar\rho$ は確率測度 $\bar\rho\,\nu$ の密度であり，
  \[
    H_\nu(\bar\rho\,\nu)
    \leq\int_\Theta H_\nu\bigl(\rho(\theta,\cdot)\,\nu\bigr)\,\d P(\theta).
  \]
\end{lemma}

\begin{proof}
  Fubini--Tonelli の定理（定理~\ref{thm:meas-product}）より
  $\bar\rho$ は可測で，
  $\int\bar\rho\,\d\nu=\int_\Theta\int_{\mathcal{X}}\rho\,\d\nu\,\d P=1$
  である．特に $\nu(\{\bar\rho=+\infty\})=0$ だから，
  $\{\bar\rho=+\infty\}$ 上で $\bar\rho$ を $0$ と置き直しても
  $\bar\rho\,\nu$ と以下の積分は変わらない．
  こうして $\bar\rho$ は $[0,\infty)$ に値を取り，
  $\bar\rho\,\nu$ は確率測度である（定義~\ref{def:meas-density}）．

  まず，各 $x\in\mathcal{X}$ を固定して
  \[
    \int_\Theta\rho(\theta,x)\log\rho(\theta,x)\,\d P(\theta)
    \geq\bar\rho(x)\log\bar\rho(x)
  \]
  を示す．
  $m:=\bar\rho(x)>0$ のとき：$t\log t\geq t-1$
  （注意~\ref{rem:meas-real-facts}）を $t=u/m$ に適用して
  $m$ 倍すれば，$u\geq0$ で
  $u\log u\geq u\log m+u-m$ を得る．
  ここで $(\rho\log\rho)^-\leq(\rho-1)^-\leq1$ は $P$-可積分だから
  左辺の積分 $\int\rho\log\rho\,\d P$ は
  $\R\cup\{+\infty\}$ の元として定まり，
  右辺 $u\log m+u-m$ は $P$-可積分（$\int\rho\,\d P=m<\infty$）である．
  $u=\rho(\theta,x)$ として積分の単調性を使うと，右辺の積分は
  $m\log m+m-m=\bar\rho(x)\log\bar\rho(x)$ になる．
  $m=0$ のとき：$\rho(\cdot,x)\geq0$ の積分が $0$ だから
  $\rho(\theta,x)=0$ が $P$-a.e. で成り立ち，両辺とも $0$ である．

  次に $g(\theta,x):=\rho\log\rho-\rho+1$ とおくと，
  $t\log t\geq t-1$ より $g\geq0$ であり，
  $\int_{\mathcal{X}}(\rho(\theta,\cdot)-1)\,\d\nu=0$ から
  \[
    H_\nu\bigl(\rho(\theta,\cdot)\,\nu\bigr)
    =\int_{\mathcal{X}}g(\theta,x)\,\d\nu(x),
    \qquad
    H_\nu(\bar\rho\,\nu)
    =\int_{\mathcal{X}}\bigl(\bar\rho\log\bar\rho-\bar\rho+1\bigr)\d\nu
  \]
  である．各 $x$ で $\rho\log\rho=g+(\rho-1)$ かつ
  $\rho-1$ は $P$-可積分だから
  $\int_\Theta g\,\d P=\int_\Theta\rho\log\rho\,\d P
  -(\bar\rho-1)$
  （$\int g\,\d P<\infty$ なら通常の線形性，
  $=+\infty$ なら両辺とも $+\infty$）であり，
  冒頭の各点評価は
  $\int_\Theta g(\theta,x)\,\d P
  \geq\bar\rho(x)\log\bar\rho(x)-\bar\rho(x)+1$
  と言い換えられる．非負可測な $g$ に Fubini--Tonelli の定理を
  適用して積分順序を交換すれば
  \[
    \int_\Theta H_\nu\bigl(\rho(\theta,\cdot)\,\nu\bigr)\,\d P
    =\int_{\mathcal{X}}\int_\Theta g\,\d P\,\d\nu
    \geq\int_{\mathcal{X}}\bigl(\bar\rho\log\bar\rho-\bar\rho+1\bigr)\d\nu
    =H_\nu(\bar\rho\,\nu).  \qedhere
  \]
\end{proof}

\begin{corollary}{ELBO の正則化項は $W_2$ を制御する}{lt-elbo-w2}
  エンコーダが密度で与えられるとする：すなわち可測な
  $\rho\colon\mathcal{X}\times Z\to[0,\infty)$ で，各 $x$ において
  $Q(x,\cdot)=\rho(x,\cdot)\,\gamma^{(d)}$ が確率測度になるとする
  （$x\mapsto Q(x,B)=\int_B\rho(x,z)\,\d\gamma^{(d)}(z)$ の可測性は
  Fubini--Tonelli の定理から従うので，$Q$ は Markov 核である）．
  このとき $q_Z=\bar\rho\,\gamma^{(d)}$
  （$\bar\rho(z)=\int\rho(x,z)\,\d P(x)$）であり，
  $q_Z\in\Pp{2}(\R^d)$ ならば
  \[
    \Wass_2\bigl(q_Z,\gamma^{(d)}\bigr)
    \leq\sqrt{2\int_{\mathcal{X}}
    H_{\gamma^{(d)}}\bigl(Q(x,\cdot)\bigr)\,\d P(x)}.
  \]
\end{corollary}

\begin{proof}
  $q_Z=\bar\rho\,\gamma^{(d)}$ は，任意の $B\in\Borel(\R^d)$ に対し
  混合測度の定義と Fubini--Tonelli の定理より
  \[
    q_Z(B)=\int_{\mathcal{X}}\int_B\rho(x,z)\,\d\gamma^{(d)}(z)\,\d P(x)
    =\int_B\bar\rho(z)\,\d\gamma^{(d)}(z)
  \]
  となることから従う．
  Talagrand の不等式（定理~\ref{thm:tc-talagrand}）と
  補題~\ref{lem:lt-entropy-convex}（$\Theta=\mathcal{X}$，
  $\nu=\gamma^{(d)}$）より
  \[
    \Wass_2(q_Z,\gamma^{(d)})^2
    \leq2\,H_{\gamma^{(d)}}(q_Z)
    \leq2\int_{\mathcal{X}}H_{\gamma^{(d)}}\bigl(Q(x,\cdot)\bigr)\,\d P(x).
    \qedhere
  \]
\end{proof}

\section{1次元の最適輸送}
\label{sec:lt-1d}

Roy--Hauberg のアルゴリズムは，画素ごとの1次元分布に対する
$W_1$ の閉形式に基づく．

\begin{theorem}{1次元の最適輸送（Villani (2003) Thm 2.18）}{lt-1d}
  $\mu,\nu\in\Prob(\R)$（$\R$ は Euclid 距離）の
  \textbf{分布関数}を
  $F(t):=\mu((-\infty,t])$，$G(t):=\nu((-\infty,t])$，
  その\textbf{一般化逆}を
  $F^{-1}(u):=\inf\{t\in\R\mid F(t)\geq u\}$（$0<u<1$）と定める．
  $h\colon[0,\infty)\to[0,\infty)$ を凸関数とし，コスト
  $c(x,y):=h(\abs{x-y})$ を考えると
  \[
    \inf_{\pi\in\Couplings(\mu,\nu)}\int c\,\d\pi
    =\int_{(0,1)}h\bigl(\abs{F^{-1}(u)-G^{-1}(u)}\bigr)\,\d\lambda_1(u).
  \]
  特に $1\leq p<\infty$，$\mu,\nu\in\Pp{p}(\R)$ に対し
  \[
    \Wass_p(\mu,\nu)^p
    =\int_{(0,1)}\abs{F^{-1}(u)-G^{-1}(u)}^p\,\d\lambda_1(u),
    \qquad
    \Wass_1(\mu,\nu)=\int_\R\abs{F(t)-G(t)}\,\d\lambda_1(t)
  \]
  であり，$F$ が連続なら単調写像 $T:=G^{-1}\circ F$ による
  像測度から定まる coupling が最適である．
  （証明は省略する．Villani,
  \emph{Topics in Optimal Transportation},
  Graduate Studies in Mathematics 58, AMS (2003),
  Theorem 2.18 および §2.2 参照．）
\end{theorem}

\section{直積測度とコストの分離}
\label{sec:lt-product}

画像は画素の直積として扱われる．
$(\mathcal{X}_i,d_i)$（$i=1,\dots,D$）を Polish 空間とし，
直積 $\mathcal{X}:=\prod_{i=1}^D\mathcal{X}_i$ に
\[
  d_p(x,y):=\Bigl(\sum_{i=1}^Dd_i(x_i,y_i)^p\Bigr)^{1/p}
  \qquad(1\leq p<\infty)
\]
を入れる．$d_p$ の三角不等式は，座標ごとの三角不等式と
$t\mapsto t^p$ の単調性（注意~\ref{rem:meas-real-facts}），
および有限和の Minkowski の不等式
（系~\ref{cor:meas-minkowski-finite}）による．
また
$\max_id_i(x_i,y_i)\leq d_p(x,y)\leq D^{1/p}\max_id_i(x_i,y_i)$
より $d_p$ は付録の $\max$ 距離と同じ開集合・同じ Cauchy 列を定め，
$(\mathcal{X},d_p)$ は Polish 空間で，$\Borel(\mathcal{X})$ は
直積 Borel $\sigma$-代数（注意~\ref{rem:meas-borel-product}）と
一致する．

\begin{theorem}{直積測度の $W_p$ は座標ごとに分解する}{lt-product}
  $\mu_i,\nu_i\in\Pp{p}(\mathcal{X}_i)$（$i=1,\dots,D$）とすると，
  $(\mathcal{X},d_p)$ 上で
  $\mu:=\mu_1\otimes\cdots\otimes\mu_D$ と
  $\nu:=\nu_1\otimes\cdots\otimes\nu_D$ は
  $\Pp{p}(\mathcal{X})$ の元であり，
  \[
    \Wass_p(\mu,\nu)^p=\sum_{i=1}^D\Wass_p(\mu_i,\nu_i)^p.
  \]
\end{theorem}

\begin{proof}
  基点 $x^0=(x^0_i)_i$ を固定する．直積測度の座標周辺分布は
  $\mu\circ\mathrm{pr}_i^{-1}=\mu_i$（矩形集合上で直接確かめられ，
  一致定理~\ref{thm:meas-uniqueness}による）だから，
  変数変換公式（補題~\ref{lem:meas-change-of-variables}）より
  \[
    \int_{\mathcal{X}}d_p(x,x^0)^p\,\d\mu
    =\sum_{i=1}^D\int_{\mathcal{X}}d_i(x_i,x^0_i)^p\,\d\mu
    =\sum_{i=1}^D\int_{\mathcal{X}_i}d_i(\cdot,x^0_i)^p\,\d\mu_i<\infty,
  \]
  すなわち $\mu\in\Pp{p}(\mathcal{X})$ である（$\nu$ も同様）．

  \textbf{（$\geq$）}
  $\pi\in\Couplings(\mu,\nu)$ を任意に取り，
  $\mathrm{pr}_{(i)}(x,y):=(x_i,y_i)$ による像測度
  $\pi_i:=\pi\circ\mathrm{pr}_{(i)}^{-1}$ を考える．
  $\mathrm{pr}_1\circ\mathrm{pr}_{(i)}=\mathrm{pr}_i\circ\mathrm{pr}_1$
  と合成則（主張~\ref{clm:meas-pf-compose}）より
  $\pi_i$ の第1周辺分布は
  $(\pi\circ\mathrm{pr}_1^{-1})\circ\mathrm{pr}_i^{-1}
  =\mu\circ\mathrm{pr}_i^{-1}=\mu_i$，
  同様に第2周辺分布は $\nu_i$ だから
  $\pi_i\in\Couplings(\mu_i,\nu_i)$ である．よって
  \[
    \int_{\mathcal{X}\times\mathcal{X}}d_p(x,y)^p\,\d\pi
    =\sum_{i=1}^D\int d_i(x_i,y_i)^p\,\d\pi
    =\sum_{i=1}^D\int d_i^p\,\d\pi_i
    \geq\sum_{i=1}^D\Wass_p(\mu_i,\nu_i)^p
  \]
  であり，$\pi$ について下限を取ればよい．

  \textbf{（$\leq$）}
  各 $i$ で最適 coupling
  $\pi_i^*\in\Couplings(\mu_i,\nu_i)$ を取る
  （定理~\ref{thm:wp-optimal-coupling-existence}）．
  直積 $\Pi:=\pi_1^*\otimes\cdots\otimes\pi_D^*$
  （$\prod_i(\mathcal{X}_i\times\mathcal{X}_i)$ 上）と，
  座標の並べ替え
  $\sigma\colon\prod_i(\mathcal{X}_i\times\mathcal{X}_i)
  \to\mathcal{X}\times\mathcal{X}$（連続ゆえ可測）による像測度
  $\tilde\pi:=\Pi\circ\sigma^{-1}$ をおく．
  矩形集合上で
  \[
    \tilde\pi\Bigl(\prod_iA_i\times\mathcal{X}\Bigr)
    =\Pi\Bigl(\prod_i(A_i\times\mathcal{X}_i)\Bigr)
    =\prod_i\pi_i^*(A_i\times\mathcal{X}_i)
    =\prod_i\mu_i(A_i)
    =\mu\Bigl(\prod_iA_i\Bigr)
  \]
  であり，矩形集合は $\pi$-系として $\Borel(\mathcal{X})$ を
  生成するから，一致定理より $\tilde\pi$ の第1周辺分布は $\mu$，
  同様に第2周辺分布は $\nu$，すなわち
  $\tilde\pi\in\Couplings(\mu,\nu)$ である．
  また $\tilde\pi$ の $(x_i,y_i)$-周辺分布は，同じ矩形集合の
  計算により $\pi_i^*$ だから，
  \[
    \Wass_p(\mu,\nu)^p
    \leq\int d_p(x,y)^p\,\d\tilde\pi
    =\sum_{i=1}^D\int d_i^p\,\d\pi_i^*
    =\sum_{i=1}^D\Wass_p(\mu_i,\nu_i)^p.  \qedhere
  \]
\end{proof}

\begin{remark}{Roy--Hauberg の式 (6) の精密化}{lt-product-audit}
  同論文の式 (6) は，等式
  $W_1^2(\otimes_i\mu_i,\otimes_i\nu_i)=\sum_iW_1^2(\mu_i,\nu_i)$
  が成り立つと主張するが，これは一般には成立しない．
  定理~\ref{thm:lt-product}が与える正確な等式は
  「ground metric の合成の指数と Wasserstein 距離の次数を
  揃えた場合」であり：
  $p=1$（$\ell^1$ 合成）なら
  $\Wass_1(\mu,\nu)=\sum_i\Wass_1(\mu_i,\nu_i)$，
  $p=2$（$\ell^2$ 合成）なら
  $\Wass_2(\mu,\nu)^2=\sum_i\Wass_2(\mu_i,\nu_i)^2$ である．
  画素ごとの $W_1$ の2乗和
  $\bigl(\sum_i\Wass_1(\mu_i,\nu_i)^2\bigr)^{1/2}$ は，
  各 $\Wass_1$ が $\Pp{1}(\mathcal{X}_i)$ 上の距離である
  （命題~\ref{prop:wp-metric}）ことと有限和の Minkowski の不等式
  （系~\ref{cor:meas-minkowski-finite}）から
  直積上の距離にはなるが，
  単一の Wasserstein 距離 $\Wass_p(\mu,\nu)$ とは
  一般に一致しない．
\end{remark}

\section{引き戻し擬距離と曲線の長さ}
\label{sec:lt-pullback}

\begin{definition}{擬距離}{lt-pseudometric}
  集合 $M$ 上の関数 $\rho\colon M\times M\to[0,\infty)$ が
  対称性・三角不等式（定義~\ref{def:meas-metric-space}の (2)(3)）と
  $\rho(z,z)=0$（$\forall z\in M$）をみたすとき，
  $\rho$ を $M$ 上の\textbf{擬距離}という
  （距離との違いは，$z\neq z'$ でも $\rho(z,z')=0$ があり得る点である）．
\end{definition}

\begin{definition}{デコーダによる引き戻し擬距離}{lt-pullback}
  $1\leq p<\infty$ とし，デコーダ $K$ は各 $z\in Z$ で
  $K(z,\cdot)\in\Pp{p}(\mathcal{X})$ をみたすとする
  （$\mathcal{X}$ は Polish 空間）．
  \[
    \rho_K(z,z'):=\Wass_p\bigl(K(z,\cdot),\,K(z',\cdot)\bigr)
    \qquad(z,z'\in Z)
  \]
  を $K$ による\textbf{引き戻し擬距離}という．
\end{definition}

\begin{claim}{$\rho_K$ は擬距離である}{lt-pullback-pseudo}
  $\rho_K$ は非負・対称で三角不等式をみたし，
  $\rho_K(z,z)=0$ である．さらに
  $z\neq z'\implies K(z,\cdot)\neq K(z',\cdot)$（$K$ が単射）ならば，
  $\rho_K$ は $Z$ 上の距離である．
\end{claim}

\begin{proof}
  $\Wass_p$ は $\Pp{p}(\mathcal{X})$ 上の距離だから
  （命題~\ref{prop:wp-metric}），すべて直ちに従う．
  一致の公理は，$\rho_K(z,z')=0\iff K(z,\cdot)=K(z',\cdot)$ による．
\end{proof}

\begin{definition}{曲線の長さと内在擬距離}{lt-length}
  $(M,\rho)$ を擬距離空間とする．
  写像 $\gamma\colon[0,1]\to M$ が \textbf{$\rho$-連続}であるとは，
  各 $t\in[0,1]$ で $s\to t$ のとき
  $\rho(\gamma(s),\gamma(t))\to0$ となることをいう．
  $\rho$-連続な $\gamma$ と
  分割 $\Delta\colon0=t_0<t_1<\cdots<t_K=1$ に対し
  \[
    S(\gamma;\Delta):=\sum_{k=1}^K
    \rho\bigl(\gamma(t_{k-1}),\gamma(t_k)\bigr),
    \qquad
    L(\gamma):=\sup_\Delta S(\gamma;\Delta)\in[0,+\infty]
  \]
  とおき，$L(\gamma)$ を $\gamma$ の\textbf{長さ}という．また
  \[
    \hat\rho(z,z')
    :=\inf\bigl\{L(\gamma)\bigm|
    \gamma\colon[0,1]\to M\text{ は }\rho\text{-連続},\
    \gamma(0)=z,\ \gamma(1)=z'\bigr\}
  \]
  を $\rho$ の定める\textbf{内在擬距離}という
  （結ぶ曲線がないときは $\inf\emptyset=+\infty$：
  定義~\ref{def:meas-sup-inf}）．
\end{definition}

\begin{claim}{長さの基本性質}{lt-length-props}
  \begin{enumerate}
    \item 分割 $\Delta'$ が $\Delta$ の細分ならば
      $S(\gamma;\Delta)\leq S(\gamma;\Delta')$ である．
    \item $L(\gamma)\geq\rho\bigl(\gamma(0),\gamma(1)\bigr)$，
      したがって $\hat\rho\geq\rho$ である．
  \end{enumerate}
\end{claim}

\begin{proof}
  (1)：分点を1つ挿入したとき，三角不等式
  $\rho(\gamma(s),\gamma(t))
  \leq\rho(\gamma(s),\gamma(r))+\rho(\gamma(r),\gamma(t))$ より
  和は減らない．細分は分点の有限回の挿入だから従う．
  (2)：自明な分割 $\Delta=\{0,1\}$ に対し
  $S(\gamma;\Delta)=\rho(\gamma(0),\gamma(1))\leq L(\gamma)$ であり，
  曲線について下限を取れば $\hat\rho\geq\rho$ である．
\end{proof}

\begin{claim}{内在擬距離は三角不等式をみたす}{lt-intrinsic}
  $\hat\rho$ は対称で，$\hat\rho(z,z)=0$，かつ
  \[
    \hat\rho(z,z'')\leq\hat\rho(z,z')+\hat\rho(z',z'')
    \qquad(\forall z,z',z''\in M)
  \]
  をみたす（値は $[0,+\infty]$ に取る）．
\end{claim}

\begin{proof}
  定数曲線は $\rho$-連続で長さ $0$ だから $\hat\rho(z,z)=0$，
  時間反転 $t\mapsto\gamma(1-t)$ は $\rho$-連続性と分割和を保つ
  （$\rho$ の対称性）から $\hat\rho$ は対称である．
  三角不等式を示す．右辺が $+\infty$ なら明らかなので，
  $\rho$-連続な $\gamma_1$（$z$ から $z'$）と
  $\gamma_2$（$z'$ から $z''$）を任意に取る．連結
  \[
    \gamma(t):=
    \begin{cases}
      \gamma_1(2t)&(0\leq t\leq1/2),\\
      \gamma_2(2t-1)&(1/2\leq t\leq1)
    \end{cases}
  \]
  は $\rho$-連続である（$t=1/2$ では
  $\gamma_1(1)=\gamma_2(0)=z'$ による）．
  任意の分割 $\Delta$ に対し，$1/2$ を加えた細分の和は
  主張~\ref{clm:lt-length-props} (1) より $S(\gamma;\Delta)$ 以上で，
  $1/2$ で二分すると $\gamma_1,\gamma_2$ の分割和の和になるから
  $S(\gamma;\Delta)\leq L(\gamma_1)+L(\gamma_2)$ である．
  分割について上限を取って
  $L(\gamma)\leq L(\gamma_1)+L(\gamma_2)$，
  さらに $\gamma_1,\gamma_2$ について下限を取ればよい．
\end{proof}

Roy--Hauberg の式 (7) は，$\rho=\rho_K$（$p=1$）に対する一様分割
$\Delta_K\colon t_k=k/K$ に沿った和 $S(\gamma;\Delta_K)$ の
$K\to\infty$ での極限として潜在空間の曲線の「長さ」を定め，
その最小化元（実装上は曲線を3次スプラインでパラメータ付けし，
曲線エネルギーを数値最適化する）を\textbf{測地線}とよぶ．
長さの定義（分割和の上限）より，式 (7) の和はつねに
$L(\gamma)$ 以下である．

\begin{remark}{文献との対応と展望}{lt-references}
  本章は Roy--Hauberg（2022，前掲）の数学的基盤の定式化である．
  定義~\ref{def:lt-vae}と系~\ref{cor:lt-elbo-w2}は同論文 §2 の
  VAE と ELBO（式 (1)）に，定理~\ref{thm:lt-1d}は式 (5) に，
  定理~\ref{thm:lt-product}と注意~\ref{rem:lt-product-audit}は
  式 (6) に，引き戻し擬距離と曲線の長さの節は式 (7) に対応する
  （式 (2)(3) の Fisher--Rao/KL 幾何は本資料の範囲外である）．
  なお同論文の prior は一般の $p(z)$ だが，本章は
  系~\ref{cor:lt-elbo-w2}で用いるため標準 Gaussian に固定した．
  集約事後分布と系~\ref{cor:lt-elbo-w2}
  （第\ref{ch:concentration}章の Talagrand の不等式による）は
  同論文にはない本章の補足である．

  残る整備項目：
  (a) 一様分割の和の極限が $L(\gamma)$ に一致すること
  （連続曲線に対する標準的事実だが，一様連続性の道具を要する），
  (b) 長さ最小化元（測地線）の存在，
  (c) 引き戻し「計量テンソル」の厳密な意味づけ —
  これは $(\Pp{2},\Wass_2)$ の Riemann 的構造
  （Benamou--Brenier の動的定式化，Otto の計算）を要し，
  測地線と曲率を扱う次の段階の主題である．
  なお同論文が比較対象とする Fisher--Rao 計量の引き戻し
  （Arvanitidis ほか 2022）は情報幾何に属し，本資料の範囲外である．
\end{remark}
